\chapter{市场环境与五个月翻倍的算术}

\section{指数稳定掩盖了高波动个股环境}

2026年7月22日，上证指数报3,867.03点，上涨0.07\%；深证成指下跌1.42\%，创业板指下跌3.23\%。两市成交约2.65万亿元，但超过3,800只个股下跌。前一交易日创业板刚刚上涨7.05\%。这组数据说明市场并不缺流动性，却缺少稳定的风险偏好。

对翻倍筛选而言，这种环境具有双重意义。一方面，短期大幅回撤把部分盈利修复标的压到较低隐含倍数；另一方面，快速轮动降低了估值重估持续到年末的概率。低价本身不是安全边际，只有现金流与盈利证据可以把它转化为安全边际。

\begin{exhibitbox}[Exhibit 2：近十个交易日回撤审计]
\small
\begin{tabularx}{\exhibitboxwidth}{L{2.3cm} R{1.5cm} R{1.5cm} R{1.5cm} R{1.7cm} X}
\toprule
标的 & 7月9日收盘 & 7月22日收盘 & 区间涨跌 & 峰谷回撤 & 估值含义 \\
\midrule
中国船舶 & 35.695 & 33.02 & -7.49\% & -12.82\% & 相对防御，但仍需重估 \\
江波龙 & 620.00 & 388.45 & -37.35\% & -39.28\% & 高拥挤、高周期定价快速出清 \\
恩捷股份 & 59.33 & 47.84 & -19.37\% & -19.37\% & 市场显著折价2027恢复 \\
盛新锂能 & 35.55 & 27.65 & -22.22\% & -26.89\% & 锂价与风险偏好双重敏感 \\
天华新能 & 69.10 & 55.88 & -19.13\% & -22.82\% & 周期重估降温 \\
雅化集团 & 20.20 & 16.79 & -16.88\% & -19.36\% & 低倍数与盈利持续性怀疑并存 \\
\bottomrule
\end{tabularx}
\sourcenote{腾讯前复权日线；2026年7月9日至7月22日。中国船舶7月20日除息0.365元，价格比较并非总回报。}
\end{exhibitbox}

\section{翻倍所需的月度复合回报}

若从7月22日持有到12月31日，约五个多月内实现100\%回报，所需月复合涨幅约为：
\[
  r_m = 2^{1/5}-1 \approx 14.9\%.
\]
这不是普通的估值修复，而是“盈利兑现、预期上修、风险偏好改善”三者共振。只有一个环节发生，通常只能带来20\%--60\%的修复，而不是翻倍。

\section{当前价格隐含的怀疑}

雅化集团现价对应2026E EPS约6.38倍，翻倍到33.58元需要12.77倍；恩捷股份现价对应2027E EPS约9.27倍，翻倍到95.68元需要18.54倍。两者并非只要盈利不下修就能翻倍，还需要倍数显著上升。

中国船舶当前约对应2027E EPS 10.55倍，翻倍需要21.10倍。江波龙当前约对应2026E EPS 14.68倍，翻倍需要29.35倍，而券商预测2027年利润增长只有约2\%。这说明江波龙的翻倍逻辑更依赖估值而不是盈利，风险显著高于表面低PE。

\begin{keyinsight}[投资含义]
近期回撤让“目标价/现价”比值迅速上升，但分子没有同步变得更可信。筛选必须同时更新现价、盈利分母、目标期限和现金流，不能把旧目标价除以新低价后直接宣布机会。
\end{keyinsight}

\clearpage
